\documentclass[9pt,t]{beamer}
% =====================================================================
% finance-beamer.tex — shared Beamer style for the LLM-in-Finance course
% =====================================================================
% Reusable slide style. Each deck \input's this immediately after
% \documentclass{beamer}. It provides:
%   * the Madrid theme + book colour palette (matches the printed book)
%   * two book-style framing blocks:
%       \begin{bigpicture} ... \end{bigpicture}   ("the bigger picture")
%       \begin{underhood}   ... \end{underhood}    ("under the hood")
%     These mirror the book's `context` and `deepdive` tcolorboxes so the
%     same intuition-vs-mechanism scaffold the reader meets in the book
%     reappears on the slides.
%   * a Python listings style for code frames
%   * appendix conventions: \appendixstart drops a divider and switches
%     to non-numbered backup frames so "extra / less central" material
%     lives after the main talk without inflating the slide count.
%
% Usage:
%   \documentclass{beamer}
%   % =====================================================================
% finance-beamer.tex — shared Beamer style for the LLM-in-Finance course
% =====================================================================
% Reusable slide style. Each deck \input's this immediately after
% \documentclass{beamer}. It provides:
%   * the Madrid theme + book colour palette (matches the printed book)
%   * two book-style framing blocks:
%       \begin{bigpicture} ... \end{bigpicture}   ("the bigger picture")
%       \begin{underhood}   ... \end{underhood}    ("under the hood")
%     These mirror the book's `context` and `deepdive` tcolorboxes so the
%     same intuition-vs-mechanism scaffold the reader meets in the book
%     reappears on the slides.
%   * a Python listings style for code frames
%   * appendix conventions: \appendixstart drops a divider and switches
%     to non-numbered backup frames so "extra / less central" material
%     lives after the main talk without inflating the slide count.
%
% Usage:
%   \documentclass{beamer}
%   % =====================================================================
% finance-beamer.tex — shared Beamer style for the LLM-in-Finance course
% =====================================================================
% Reusable slide style. Each deck \input's this immediately after
% \documentclass{beamer}. It provides:
%   * the Madrid theme + book colour palette (matches the printed book)
%   * two book-style framing blocks:
%       \begin{bigpicture} ... \end{bigpicture}   ("the bigger picture")
%       \begin{underhood}   ... \end{underhood}    ("under the hood")
%     These mirror the book's `context` and `deepdive` tcolorboxes so the
%     same intuition-vs-mechanism scaffold the reader meets in the book
%     reappears on the slides.
%   * a Python listings style for code frames
%   * appendix conventions: \appendixstart drops a divider and switches
%     to non-numbered backup frames so "extra / less central" material
%     lives after the main talk without inflating the slide count.
%
% Usage:
%   \documentclass{beamer}
%   \input{../_style/finance-beamer.tex}
%   \title{...}\subtitle{...}\author{...}\institute{...}\date{\today}
%   \begin{document} ... \end{document}
% =====================================================================

\usetheme{Madrid}
\usecolortheme{whale}
\setbeamertemplate{navigation symbols}{}
\setbeamertemplate{caption}[numbered]
\setbeamercovered{transparent}

% --- Density controls: keep dense, equation-heavy frames on one slide ---
% Slightly smaller body text + tighter lists/blocks. Frame titles and the
% Madrid chrome keep their normal size, so the deck still reads as Madrid,
% just with more vertical room for content.
\setbeamerfont{normal text}{size=\small}
\setbeamertemplate{itemize/enumerate body begin}{\small}
\setbeamertemplate{itemize/enumerate subbody begin}{\footnotesize}
% Tighter block spacing.
\setbeamertemplate{block begin}{%
  \par\vskip2pt%
  \begin{beamercolorbox}[colsep*=.75ex,leftskip=1ex,rightskip=1ex]{block title}%
    \usebeamerfont*{block title}\insertblocktitle%
  \end{beamercolorbox}%
  {\parskip0pt\vskip-1pt}%
  \begin{beamercolorbox}[colsep*=.75ex,vmode,leftskip=1ex,rightskip=1ex]{block body}%
  \vskip1pt}
\setbeamertemplate{block end}{\end{beamercolorbox}\vskip2pt}

\usepackage{amsmath, amssymb}
\usepackage{booktabs}
\usepackage{xcolor}
\usepackage{listings}
\usepackage[skins, breakable]{tcolorbox}

% --- Book colour palette (kept in sync with book/preamble.tex) ---
\definecolor{linkblue}{RGB}{14, 75, 140}
\definecolor{deepdivebg}{RGB}{235, 244, 255}
\definecolor{narrativebg}{RGB}{255, 252, 240}
\definecolor{narrativerule}{RGB}{180, 130, 40}
\definecolor{steelblue}{RGB}{70, 130, 180}
\definecolor{forestgreen}{RGB}{34, 139, 34}
\definecolor{crimson}{RGB}{220, 20, 60}
\definecolor{darkorange}{RGB}{255, 140, 0}
\definecolor{codebg}{RGB}{248, 248, 248}
\definecolor{coderule}{RGB}{210, 210, 210}
\definecolor{keywordblue}{RGB}{0, 100, 180}
\definecolor{commentgray}{RGB}{120, 120, 120}
\definecolor{stringgreen}{RGB}{30, 130, 30}

% Tie the Madrid structural colour to the book's link blue.
\setbeamercolor{structure}{fg=linkblue}
\setbeamercolor{title}{fg=white}
\setbeamercolor{frametitle}{fg=white}

% --- "the bigger picture" : narrative / intuition framing -----------
\newtcolorbox{bigpicture}{%
  enhanced, breakable, boxsep=2pt,
  colback  = narrativebg,
  colframe = narrativerule,
  boxrule  = 0pt, toprule = 1.4pt, bottomrule = 0.4pt,
  leftrule = 0pt, rightrule = 0pt, arc = 0pt,
  left = 6pt, right = 6pt, top = 4pt, bottom = 5pt,
  fonttitle = \footnotesize\scshape\color{narrativerule},
  title = {the bigger picture}, attach title to upper = {\par\smallskip},
}

% --- "under the hood" : the mechanism / technical detail ------------
\newtcolorbox{underhood}{%
  enhanced, breakable, boxsep=2pt,
  colback  = deepdivebg,
  colframe = linkblue,
  boxrule  = 0pt, toprule = 1.4pt, bottomrule = 0.4pt,
  leftrule = 0pt, rightrule = 0pt, arc = 0pt,
  left = 6pt, right = 6pt, top = 4pt, bottom = 5pt,
  fonttitle = \footnotesize\scshape\color{linkblue},
  title = {under the hood}, attach title to upper = {\par\smallskip},
}

% --- Python code style ---------------------------------------------
\lstset{
  basicstyle    = \ttfamily\footnotesize,
  keywordstyle  = \color{keywordblue}\bfseries,
  commentstyle  = \color{commentgray}\itshape,
  stringstyle   = \color{stringgreen},
  showstringspaces = false,
  language      = Python,
  breaklines    = true,
  backgroundcolor = \color{codebg},
  frame         = single,
  rulecolor     = \color{coderule},
  numbers       = none,
  aboveskip     = 4pt, belowskip = 4pt,
}

% --- Appendix conventions ------------------------------------------
% \appendixstart{Title} : begins the "extra material" appendix. Frames
% after it are not counted toward the main deck's frame total, keeping
% the core talk lean while preserving the deeper / optional content.
\newcommand{\appendixstart}[1]{%
  \appendix
  \setbeamertemplate{frametitle continuation}{}%
  \begin{frame}[noframenumbering]
    \begin{center}
      {\Large\scshape\color{linkblue} Appendix}\\[4pt]
      {\large #1}\\[10pt]
      {\footnotesize Extra / optional material — beyond the core lecture.}
    \end{center}
  \end{frame}
}
% Use \begin{frame}[noframenumbering]{...} for each appendix frame.

%   \title{...}\subtitle{...}\author{...}\institute{...}\date{\today}
%   \begin{document} ... \end{document}
% =====================================================================

\usetheme{Madrid}
\usecolortheme{whale}
\setbeamertemplate{navigation symbols}{}
\setbeamertemplate{caption}[numbered]
\setbeamercovered{transparent}

% --- Density controls: keep dense, equation-heavy frames on one slide ---
% Slightly smaller body text + tighter lists/blocks. Frame titles and the
% Madrid chrome keep their normal size, so the deck still reads as Madrid,
% just with more vertical room for content.
\setbeamerfont{normal text}{size=\small}
\setbeamertemplate{itemize/enumerate body begin}{\small}
\setbeamertemplate{itemize/enumerate subbody begin}{\footnotesize}
% Tighter block spacing.
\setbeamertemplate{block begin}{%
  \par\vskip2pt%
  \begin{beamercolorbox}[colsep*=.75ex,leftskip=1ex,rightskip=1ex]{block title}%
    \usebeamerfont*{block title}\insertblocktitle%
  \end{beamercolorbox}%
  {\parskip0pt\vskip-1pt}%
  \begin{beamercolorbox}[colsep*=.75ex,vmode,leftskip=1ex,rightskip=1ex]{block body}%
  \vskip1pt}
\setbeamertemplate{block end}{\end{beamercolorbox}\vskip2pt}

\usepackage{amsmath, amssymb}
\usepackage{booktabs}
\usepackage{xcolor}
\usepackage{listings}
\usepackage[skins, breakable]{tcolorbox}

% --- Book colour palette (kept in sync with book/preamble.tex) ---
\definecolor{linkblue}{RGB}{14, 75, 140}
\definecolor{deepdivebg}{RGB}{235, 244, 255}
\definecolor{narrativebg}{RGB}{255, 252, 240}
\definecolor{narrativerule}{RGB}{180, 130, 40}
\definecolor{steelblue}{RGB}{70, 130, 180}
\definecolor{forestgreen}{RGB}{34, 139, 34}
\definecolor{crimson}{RGB}{220, 20, 60}
\definecolor{darkorange}{RGB}{255, 140, 0}
\definecolor{codebg}{RGB}{248, 248, 248}
\definecolor{coderule}{RGB}{210, 210, 210}
\definecolor{keywordblue}{RGB}{0, 100, 180}
\definecolor{commentgray}{RGB}{120, 120, 120}
\definecolor{stringgreen}{RGB}{30, 130, 30}

% Tie the Madrid structural colour to the book's link blue.
\setbeamercolor{structure}{fg=linkblue}
\setbeamercolor{title}{fg=white}
\setbeamercolor{frametitle}{fg=white}

% --- "the bigger picture" : narrative / intuition framing -----------
\newtcolorbox{bigpicture}{%
  enhanced, breakable, boxsep=2pt,
  colback  = narrativebg,
  colframe = narrativerule,
  boxrule  = 0pt, toprule = 1.4pt, bottomrule = 0.4pt,
  leftrule = 0pt, rightrule = 0pt, arc = 0pt,
  left = 6pt, right = 6pt, top = 4pt, bottom = 5pt,
  fonttitle = \footnotesize\scshape\color{narrativerule},
  title = {the bigger picture}, attach title to upper = {\par\smallskip},
}

% --- "under the hood" : the mechanism / technical detail ------------
\newtcolorbox{underhood}{%
  enhanced, breakable, boxsep=2pt,
  colback  = deepdivebg,
  colframe = linkblue,
  boxrule  = 0pt, toprule = 1.4pt, bottomrule = 0.4pt,
  leftrule = 0pt, rightrule = 0pt, arc = 0pt,
  left = 6pt, right = 6pt, top = 4pt, bottom = 5pt,
  fonttitle = \footnotesize\scshape\color{linkblue},
  title = {under the hood}, attach title to upper = {\par\smallskip},
}

% --- Python code style ---------------------------------------------
\lstset{
  basicstyle    = \ttfamily\footnotesize,
  keywordstyle  = \color{keywordblue}\bfseries,
  commentstyle  = \color{commentgray}\itshape,
  stringstyle   = \color{stringgreen},
  showstringspaces = false,
  language      = Python,
  breaklines    = true,
  backgroundcolor = \color{codebg},
  frame         = single,
  rulecolor     = \color{coderule},
  numbers       = none,
  aboveskip     = 4pt, belowskip = 4pt,
}

% --- Appendix conventions ------------------------------------------
% \appendixstart{Title} : begins the "extra material" appendix. Frames
% after it are not counted toward the main deck's frame total, keeping
% the core talk lean while preserving the deeper / optional content.
\newcommand{\appendixstart}[1]{%
  \appendix
  \setbeamertemplate{frametitle continuation}{}%
  \begin{frame}[noframenumbering]
    \begin{center}
      {\Large\scshape\color{linkblue} Appendix}\\[4pt]
      {\large #1}\\[10pt]
      {\footnotesize Extra / optional material — beyond the core lecture.}
    \end{center}
  \end{frame}
}
% Use \begin{frame}[noframenumbering]{...} for each appendix frame.

%   \title{...}\subtitle{...}\author{...}\institute{...}\date{\today}
%   \begin{document} ... \end{document}
% =====================================================================

\usetheme{Madrid}
\usecolortheme{whale}
\setbeamertemplate{navigation symbols}{}
\setbeamertemplate{caption}[numbered]
\setbeamercovered{transparent}

% --- Density controls: keep dense, equation-heavy frames on one slide ---
% Slightly smaller body text + tighter lists/blocks. Frame titles and the
% Madrid chrome keep their normal size, so the deck still reads as Madrid,
% just with more vertical room for content.
\setbeamerfont{normal text}{size=\small}
\setbeamertemplate{itemize/enumerate body begin}{\small}
\setbeamertemplate{itemize/enumerate subbody begin}{\footnotesize}
% Tighter block spacing.
\setbeamertemplate{block begin}{%
  \par\vskip2pt%
  \begin{beamercolorbox}[colsep*=.75ex,leftskip=1ex,rightskip=1ex]{block title}%
    \usebeamerfont*{block title}\insertblocktitle%
  \end{beamercolorbox}%
  {\parskip0pt\vskip-1pt}%
  \begin{beamercolorbox}[colsep*=.75ex,vmode,leftskip=1ex,rightskip=1ex]{block body}%
  \vskip1pt}
\setbeamertemplate{block end}{\end{beamercolorbox}\vskip2pt}

\usepackage{amsmath, amssymb}
\usepackage{booktabs}
\usepackage{xcolor}
\usepackage{listings}
\usepackage[skins, breakable]{tcolorbox}

% --- Book colour palette (kept in sync with book/preamble.tex) ---
\definecolor{linkblue}{RGB}{14, 75, 140}
\definecolor{deepdivebg}{RGB}{235, 244, 255}
\definecolor{narrativebg}{RGB}{255, 252, 240}
\definecolor{narrativerule}{RGB}{180, 130, 40}
\definecolor{steelblue}{RGB}{70, 130, 180}
\definecolor{forestgreen}{RGB}{34, 139, 34}
\definecolor{crimson}{RGB}{220, 20, 60}
\definecolor{darkorange}{RGB}{255, 140, 0}
\definecolor{codebg}{RGB}{248, 248, 248}
\definecolor{coderule}{RGB}{210, 210, 210}
\definecolor{keywordblue}{RGB}{0, 100, 180}
\definecolor{commentgray}{RGB}{120, 120, 120}
\definecolor{stringgreen}{RGB}{30, 130, 30}

% Tie the Madrid structural colour to the book's link blue.
\setbeamercolor{structure}{fg=linkblue}
\setbeamercolor{title}{fg=white}
\setbeamercolor{frametitle}{fg=white}

% --- "the bigger picture" : narrative / intuition framing -----------
\newtcolorbox{bigpicture}{%
  enhanced, breakable, boxsep=2pt,
  colback  = narrativebg,
  colframe = narrativerule,
  boxrule  = 0pt, toprule = 1.4pt, bottomrule = 0.4pt,
  leftrule = 0pt, rightrule = 0pt, arc = 0pt,
  left = 6pt, right = 6pt, top = 4pt, bottom = 5pt,
  fonttitle = \footnotesize\scshape\color{narrativerule},
  title = {the bigger picture}, attach title to upper = {\par\smallskip},
}

% --- "under the hood" : the mechanism / technical detail ------------
\newtcolorbox{underhood}{%
  enhanced, breakable, boxsep=2pt,
  colback  = deepdivebg,
  colframe = linkblue,
  boxrule  = 0pt, toprule = 1.4pt, bottomrule = 0.4pt,
  leftrule = 0pt, rightrule = 0pt, arc = 0pt,
  left = 6pt, right = 6pt, top = 4pt, bottom = 5pt,
  fonttitle = \footnotesize\scshape\color{linkblue},
  title = {under the hood}, attach title to upper = {\par\smallskip},
}

% --- Python code style ---------------------------------------------
\lstset{
  basicstyle    = \ttfamily\footnotesize,
  keywordstyle  = \color{keywordblue}\bfseries,
  commentstyle  = \color{commentgray}\itshape,
  stringstyle   = \color{stringgreen},
  showstringspaces = false,
  language      = Python,
  breaklines    = true,
  backgroundcolor = \color{codebg},
  frame         = single,
  rulecolor     = \color{coderule},
  numbers       = none,
  aboveskip     = 4pt, belowskip = 4pt,
}

% --- Appendix conventions ------------------------------------------
% \appendixstart{Title} : begins the "extra material" appendix. Frames
% after it are not counted toward the main deck's frame total, keeping
% the core talk lean while preserving the deeper / optional content.
\newcommand{\appendixstart}[1]{%
  \appendix
  \setbeamertemplate{frametitle continuation}{}%
  \begin{frame}[noframenumbering]
    \begin{center}
      {\Large\scshape\color{linkblue} Appendix}\\[4pt]
      {\large #1}\\[10pt]
      {\footnotesize Extra / optional material — beyond the core lecture.}
    \end{center}
  \end{frame}
}
% Use \begin{frame}[noframenumbering]{...} for each appendix frame.


\title{Practical Session 17: Building Agentic Iteration Loops}
\subtitle{Hands-On: A Root-Finding Loop and a Credit-Risk Stress Loop}
\author{Juan F.\ Imbet}
\institute{Paris Dauphine -- PSL University}
\date{\today}

\begin{document}

\begin{frame}\titlepage\end{frame}

\begin{frame}{Outline}\tableofcontents\end{frame}

% =====================================================================
\section{Recap: What You Need}
% =====================================================================

\begin{frame}{Recap: The Loop and Its Three Building Blocks}

\begin{block}{The goal-directed loop}
\textbf{Observe} $\to$ \textbf{Propose} $\to$ \textbf{Act} $\to$
\textbf{Evaluate} $\to$ \textbf{Revise}, until $\mathrm{ok}(s)$ holds or a
stopping rule fires. Genuine iteration: $a_{t+1}$ depends on feedback $e_{t+1}$.
\end{block}

\begin{itemize}
  \item \textbf{Agents} --- \emph{who} acts (worker, critic, validator, planner,
        refactorer). Value $=$ the durable change to state.
  \item \textbf{Skills} --- \emph{how} recurring work is done the same way twice
        (a skill $\approx$ a function).
  \item \textbf{Hooks} --- \emph{what} the work may never violate; fire
        automatically, outside the agent's control.
\end{itemize}

\medskip
\begin{alertblock}{The rule we will exploit twice today}
Never let the agent's own ``done'' end the loop. Termination must hang on a check
the agent does \emph{not} control --- a \alert{tolerance hook}.
\end{alertblock}
\end{frame}

\begin{frame}{Recap: Formulas and Rules for Today's Problems}

\textbf{Root finding (sign-and-magnitude search):} from a bracket $[a,b]$ with
$f(a)<0<f(b)$, the root lies between; weight the next guess toward the endpoint
with the \emph{smaller} $|f|$ (false-position intuition). Success: $|f(x)|<10^{-3}$.

\medskip
\textbf{Expected loss:}
\[
  \mathrm{EL} \;=\; \mathrm{PD}\times\mathrm{LGD}\times\mathrm{EAD},
  \qquad
  \mathrm{EL}\,\text{(bps of EAD)} \;=\; \mathrm{PD}\times\mathrm{LGD}\times 10^4 .
\]

\medskip
\textbf{Illustrative bank credit policy (today's success criterion):}
\begin{itemize}
  \item Stressed EL $\leq 150$ bps $\Rightarrow$ \emph{lend} (clean).
  \item $150 <$ stressed EL $\leq 500$ bps $\Rightarrow$ \emph{conditional / modify}.
  \item Stressed EL $> 500$ bps $\Rightarrow$ \emph{decline}.
  \item ``Low risk'' permitted \emph{only if} base $\leq 50$ bps \emph{and}
        stressed $\leq 150$ bps.
\end{itemize}

\medskip
\footnotesize All figures today are illustrative and internally consistent
($\mathrm{EL}=\mathrm{PD}\cdot\mathrm{LGD}\cdot\mathrm{EAD}$ in every column).
\end{frame}

% =====================================================================
\section{Guided Problems}
% =====================================================================

\begin{frame}{How to Work the Next Three Problems}

\begin{enumerate}
  \item \textbf{P1 --- IRR loop (transfer):} apply Example~1's sign-and-magnitude
        search to solve for an IRR from feedback alone.
  \item \textbf{P2 --- The success gate:} write the tolerance hook in bash that
        vetoes a premature ``success.''
  \item \textbf{P3 --- Credit mini case:} compute base vs.\ stressed EL and make
        the lend / modify / decline call.
\end{enumerate}

\medskip
\begin{bigpicture}
P1 and P2 are the \emph{same loop} from two angles: P1 is the agent's directed
search; P2 is the guardrail that defines ``done.'' P3 scales the idea up to a
judgment a committee must accept.
\end{bigpicture}

\medskip
\textbf{Companion notebooks:} \texttt{code/notebooks/17-loops-goals-iterations/}
--- run the evaluator, the stop hook, and the EL reconciliation end-to-end.
\end{frame}

\begin{frame}{P1 --- Apply the Loop to an IRR (Transfer Task)}

\textbf{This is Example~1's loop with $x$ renamed $r$ and $f$ renamed NPV ---
the same machine, only the residual changes.}

\medskip
\textbf{Project cash flows:} $-100$ at $t=0$, $+60$ at $t=1$, $+60$ at $t=2$.
Find the rate $r$ (the IRR) with
\[
  \mathrm{NPV}(r) = -100 + \frac{60}{1+r} + \frac{60}{(1+r)^2} = 0,
  \qquad \text{tolerance } |\mathrm{NPV}|<0.05 .
\]
Derivatives are \alert{forbidden} --- you may only evaluate NPV at a trial rate
and read the result. (Sign rule: $\mathrm{NPV}>0$ $\Rightarrow$ the rate is too
\emph{low}, so raise $r$.)

\medskip
The loop has already run two iterations:

\begin{table}
\centering
\begin{tabular}{@{}rrl@{}}
\toprule
$r$ & $\mathrm{NPV}(r)$ & Status \\
\midrule
$0.00$ & $+20.00$ & Positive --- rate too low \\
$0.20$ & $-8.33$ & Negative --- overshot \\
\bottomrule
\end{tabular}
\end{table}

\medskip
\textbf{Your tasks:}
\begin{enumerate}
  \item Confirm the first two evaluations \emph{bracket} the root. Where does the
        IRR lie?
  \item From the \emph{sign and magnitude} of NPV, propose the next two trial
        rates $r_3,\,r_4$.
  \item Are you within tolerance yet? Which loop component decides --- the agent or
        the hook?
\end{enumerate}
\end{frame}

\begin{frame}{P2 --- Write the Tolerance Success-Gate Hook \emph{from spec}}

\textbf{Do not copy code --- build it from the specification.} The IRR evaluator
of P1 appends one line per call to \texttt{run-log.txt}, the last line being the
current state on disk.

\medskip
\begin{block}{Specification}
Write a bash stop hook that:
\begin{itemize}
  \item reads the \emph{latest} logged NPV from \texttt{run-log.txt};
  \item exits \texttt{0} (allow ``success'') \emph{only if} $|\mathrm{NPV}|<0.05$;
  \item otherwise prints a clear, human-readable message and exits \texttt{1},
        which \alert{blocks} the agent from declaring success.
\end{itemize}
\end{block}

\medskip
\begin{alertblock}{Why this matters}
This single script is the difference between \emph{feedback-driven iteration} and
\emph{premature success}. Without it, the agent can announce ``done'' at
$r=0.135$ ($|\mathrm{NPV}|=0.56$) and be believed.
\end{alertblock}

\medskip
\footnotesize \emph{Plan first:} how do you grab the last line's number
(\texttt{tail}/\texttt{grep}/\texttt{cut}), and how do you test an absolute-value
inequality in shell? Then write it. The reference solution follows later.
\end{frame}

\begin{frame}{P3 --- Credit Mini Case: Atlas Foods}

\textbf{Atlas Foods} (fictional; figures illustrative) seeks a committed revolver.
\alert{EAD $=$ \$20\,m.} The base-case PD model and LGD agent report:

\medskip
\begin{tabular}{@{}lrr@{}}
\toprule
Input & Base case & Stressed (adverse) \\
\midrule
Probability of default (PD) & $2.0\%$ & $8.0\%$ \\
Loss given default (LGD) & $35\%$ & $50\%$ \\
Exposure at default (EAD) & \$20\,m & \$20\,m \\
\bottomrule
\end{tabular}

\medskip
The stressed PD and LGD are the loop's outputs after applying a combined shock
(revenue $-20\%$, margin $-300$ bps, rate $+250$ bps, DSO $+25$ days), with
recovery deteriorating on illiquid collateral.

\medskip
\textbf{Your tasks:}
\begin{enumerate}
  \item Compute base and stressed EL, in \emph{dollars} and in \emph{bps of EAD}.
  \item Apply the policy (recap slide): lend / modify / decline?
  \item May the memo say ``low risk''? Which \emph{hook} decides, and why is that
        better than leaving it to the writer?
\end{enumerate}
\end{frame}

% =====================================================================
\section{Mini Case Study and Discussion}
% =====================================================================

\begin{frame}{Mini Case Study: Designing the Atlas Loop}

You are asked to wrap the Atlas analysis in a \emph{goal-directed loop} rather
than a one-shot memo. Sketch, in two minutes with a neighbour:

\medskip
\begin{itemize}
  \item \textbf{Goal file} --- state the success criterion, a stopping rule, and
        one anticipated failure mode.
  \item \textbf{Agents} --- which two or three roles are genuinely needed (resist
        excessive delegation)?
  \item \textbf{Hooks} --- which automatic checks must fire, and \emph{when}
        (pre- vs.\ post-action)?
  \item \textbf{Escalation trigger} --- what condition forces a human into the
        loop before it proceeds?
  \item \textbf{Pattern selection} --- which of the lecture's \emph{five} design
        patterns best fits the Atlas credit workflow? Name it and justify it.
\end{itemize}

\medskip
\begin{bigpicture}
There is no single right answer --- but a good design names a \emph{checkable}
goal, a \emph{guaranteed-terminating} stopping rule, at least one
\emph{outcome-tied} hook, and an escalation point that does not depend on the
agent's goodwill.
\end{bigpicture}
\end{frame}

\begin{frame}{Discussion Prompt}

\begin{block}{Debate (5 minutes)}
\textbf{``The Atlas loop converged: every hook is green and the skeptic signed
off. Should the bank lend?''}
\end{block}

\medskip
Argue both sides, using the lecture's vocabulary:
\begin{itemize}
  \item \textbf{Goodhart / fake convergence:} could the hooks be green while the
        true goal is unmet? What would gaming look like here?
  \item \textbf{Accountability:} the skeptic's sign-off is \emph{not} sign-off.
        Who owns the decision, and what does the unsigned block add?
  \item \textbf{Audit trail:} if a regulator asks ``where did the stressed
        $8\%$ PD come from?'' three months later, can you answer?
\end{itemize}

\medskip
\begin{alertblock}{The point}
A green loop produces a \emph{defensible recommendation}, not a decision. ``Lend''
is a human act that the loop can evidence but never authorize.
\end{alertblock}
\end{frame}

% =====================================================================
\section{Solutions}
% =====================================================================

\begin{frame}{P1 Solution --- The IRR Bracket Tightens}

\textbf{(a)} $\mathrm{NPV}(0.00)=+20.00>0$ and $\mathrm{NPV}(0.20)=-8.33<0$: NPV
goes positive$\to$negative, so it crosses zero in between --- the IRR lies in
$(0.00,\,0.20)$. \textbf{(b)} Linear interpolation across the bracket,
$0+\tfrac{20}{20+8.33}\times 0.20\approx 0.14$, points near $14\%$; test there
and tighten on each sign change:

\begin{table}
\centering
\begin{tabular}{@{}rrl@{}}
\toprule
$r$ & $\mathrm{NPV}(r)$ & Reasoning \\
\midrule
$0.00$ & $+20.00$ & Positive; rate too low. \\
$0.20$ & $-8.33$ & Negative; bracket $(0.00,0.20)$. \\
$\mathbf{0.13}$ & $\mathbf{+0.09}$ & Slightly positive; root in $(0.13,0.20)$. \\
$\mathbf{0.135}$ & $\mathbf{-0.56}$ & Negative; root in $(0.13,0.135)$. \\
$0.131$ & $-0.04$ & $|\mathrm{NPV}|<0.05$: IRR $\approx 13.1\%$. \\
\bottomrule
\end{tabular}
\end{table}

\smallskip
\textbf{The same machine.} This is literally Example~1's sign-and-magnitude
search with $x\to r$ and $f\to\mathrm{NPV}$: bracket on a sign change, then weight
the next guess toward the endpoint with the smaller $|\mathrm{NPV}|$.

\medskip
\begin{block}{Who decides ``done''?}
\alert{The hook, not the agent.} At $r=0.135$, $|\mathrm{NPV}|=0.56$ --- still
above tolerance. The agent proposes; the tolerance hook (P2) is the only thing
allowed to certify convergence.
\end{block}
\end{frame}

\begin{frame}[fragile]{P2 Solution --- The Success-Gate Hook}

\begin{lstlisting}[language=bash]
#!/usr/bin/env bash
# stop-hook: veto "success" unless the latest |NPV| < 0.05
set -euo pipefail

last=$(grep -o 'NPV=[-0-9.eE]*' run-log.txt | tail -n1 | cut -d= -f2)
ok=$(python -c "print(1 if abs(float('$last')) < 0.05 else 0)")

if [ "$ok" -ne 1 ]; then
  echo "BLOCKED: |NPV|=$last >= 0.05 -- not converged" >&2
  exit 1          # nonzero exit = veto: the loop may NOT stop
fi
echo "OK: |NPV| < 0.05 -- converged"
exit 0
\end{lstlisting}

\smallskip
\textbf{Key properties:} it reads the \emph{last} line (the current state on
disk), it is \emph{deterministic} (same log $\to$ same verdict), and the agent
cannot skip it --- the harness treats the nonzero exit as a hard veto. This is the
seven-step template's step 5 (validate) made real.
\end{frame}

\begin{frame}{P3 Solution --- Atlas Foods Expected Loss}

\textbf{Base case:}
\[
  \mathrm{EL} = 0.020 \times 0.35 \times \$20\text{m} = \$0.14\text{m}
  \;=\; \mathbf{70}\text{ bps of EAD}.
\]
\textbf{Stressed (adverse):}
\[
  \mathrm{EL} = 0.080 \times 0.50 \times \$20\text{m} = \$0.80\text{m}
  \;=\; \mathbf{400}\text{ bps of EAD}.
\]

\medskip
\begin{block}{Decision}
Stressed EL $=400$ bps falls in $(150, 500]$ $\Rightarrow$
\alert{conditional / modify} --- not a clean lend, not a decline. Reasonable
conditions: a smaller facility, tighter covenants, more collateral (to cut LGD),
and pricing that compensates for the stressed EL.
\end{block}

\medskip
\textbf{May the memo say ``low risk''?} No. The \emph{``low risk'' gate} requires
base $\leq 50$ \emph{and} stressed $\leq 150$ bps; here base $=70$ and stressed
$=400$. \alert{The hook decides, not the writer} --- because language is a control
surface, and the phrase must be earned by the arithmetic.
\end{frame}

\begin{frame}{Mini-Case Solution --- Which Design Pattern?}

\begin{block}{Intended answer}
The \alert{credit-model: build $\to$ stress-test $\to$ red-team $\to$ finalize}
pattern --- a specialization of \emph{scaffold--draft--audit--repair} with an
explicit adversarial skeptic stage.
\end{block}

\medskip
\textbf{Why this one, not the others:}
\begin{itemize}
  \item The goal is a \emph{judgment} (an acceptable committee memo), not a single
        machine-checkable number, so \emph{generate--test--revise} alone is too
        thin.
  \item It is \emph{high-stakes}: the two stages that distinguish this pattern ---
        \emph{mandatory stress testing} (adverse and severe, not just base) and a
        \emph{red-team} skeptic that demands sensitivities and peer benchmarks ---
        are exactly what Atlas needs.
  \item It \emph{finalizes} into a governed artifact with an unsigned sign-off
        block, encoding SR~11-7 governance. Sub-patterns still nest inside (the
        PD-model build stage is generate--test--revise; the red-team is
        planner--executor--critic).
\end{itemize}
\end{frame}

\begin{frame}{Wrap-Up}

\textbf{What you practised:}
\begin{itemize}
  \item Directed search from \emph{feedback alone} (P1) --- the engine inside every
        loop in the book.
  \item A \emph{tolerance hook} (P2) that makes ``done'' external to the agent.
  \item EL reconciliation and a \emph{policy-gated} lend decision (P3) --- the
        credit loop in miniature.
\end{itemize}

\medskip
\begin{bigpicture}
The root-finder and the credit memo are the \emph{same machine}: a state variable,
a validator, and a definition of ``residual.'' Only the residual changes --- a
number, a failing test, or a committee's objection.
\end{bigpicture}

\medskip
\textbf{Take it further:} \texttt{code/notebooks/17-loops-goals-iterations/} runs
both loops end-to-end. See book Chapter 17 and Appendices A--F for the full
treatment and the concrete tooling.
\end{frame}

\end{document}
